\usepackage[a4paper,left=23mm,right=27mm,top=20mm,bottom=21mm]{geometry}
\usepackage{amsmath,amssymb,booktabs,tabularx,array}
\usepackage{graphicx,xcolor,tikz}
\usetikzlibrary{arrows.meta,positioning,calc}
\usepackage{caption,float}
\usepackage[hidelinks,unicode]{hyperref}
\usepackage{url}
\definecolor{ink}{HTML}{17364B}
\definecolor{accent}{HTML}{236B81}
\definecolor{pale}{HTML}{EEF4F6}
\definecolor{muted}{HTML}{50616A}
\ctexset{section={format=\Large\bfseries\color{ink},beforeskip=1ex,afterskip=1ex},subsection={format=\large\bfseries,beforeskip=1ex,afterskip=.7ex}}
\captionsetup{font=small,labelfont=bf,labelsep=quad,skip=5pt}
\setlength{\parindent}{2em}
\setlength{\parskip}{.36em}
\setlength{\emergencystretch}{2em}
\setlength{\tabcolsep}{4pt}
\renewcommand{\arraystretch}{1.16}
\linespread{1.17}
\pagestyle{plain}
\graphicspath{{../assets/figures/}}
\hypersetup{pdfauthor={},pdfsubject={DRAM单粒子效应仿真研究},pdfcreator={XeLaTeX}}
\newcommand{\um}{\,\mu\mathrm{m}}
\newcommand{\ns}{\,\mathrm{ns}}
\newcommand{\mV}{\,\mathrm{mV}}
\newcommand{\fC}{\,\mathrm{fC}}
\newcommand{\V}{\,\mathrm{V}}
\newcommand{\fig}[4]{\begin{figure}[H]\centering\includegraphics[width=#2\linewidth]{#1}\caption{#3}\label{#4}\end{figure}}
\newcommand{\atable}[3]{\begin{table}[H]\centering\small\caption{#2}\label{#3}\input{../assets/tables/#1.tex}\end{table}}
\newcommand{\pageheading}[1]{\clearpage\section{#1}}
\newcommand{\smallnote}[1]{\par\noindent{\small\color{muted}#1}\par}
\newcommand{\modeldiagram}{%
\begin{tikzpicture}[font=\small,box/.style={draw=accent,fill=pale,rounded corners=1pt,minimum height=1.1cm,text width=3.35cm,align=center},arr/.style={-{Latex[length=2mm]},draw=accent,thick}]
\node[box] (bl) {位线 BL\\$C_b=100\,\mathrm{fF}$};
\node[box,right=1.0cm of bl] (mos) {二维访问晶体管\\栅长 $L$};
\node[box,right=1.0cm of mos] (sn) {存储节点 SN\\$C_s=10\,\mathrm{fF}$};
\node[above=.6cm of mos] (wl) {字线 WL：控制导通};
\draw[<->,thick,draw=accent] (bl)--(mos);
\draw[<->,thick,draw=accent] (mos)--(sn);
\draw[arr] (wl)--(mos);
\node[below=.65cm of sn,align=center] {初态 $0.1/0.9\,\mathrm V$\\分别对应存储0/1};
\node[below=.65cm of bl,align=center] {读出前预充至$0.5\,\mathrm V$\\经$50\,\Omega$支路};
\end{tikzpicture}}
\newcommand{\workflowdiagram}{%
\begin{tikzpicture}[font=\small,box/.style={draw=accent,fill=pale,minimum height=1.05cm,text width=3.55cm,align=center},arr/.style={-{Latex[length=2mm]},draw=accent,thick},node distance=.7cm and .7cm]
\node[box] (request) {中文参数请求\\检查范围和单位};
\node[box,right=of request] (g4) {Geant4\\逐步电离沉积};
\node[box,right=of g4] (source) {等效电荷源\\总量与径迹};
\node[box,below=of source] (tcad) {TCAD瞬态\\信号、对照和参考};
\node[box,left=of tcad] (metrics) {电荷变化与读出裕量\\$\Delta Q,\ M,\ \Delta M$};
\node[box,left=of metrics] (score) {固定预测评分\\留出与外推};
\draw[arr] (request)--(g4);\draw[arr] (g4)--(source);\draw[arr] (source)--(tcad);\draw[arr] (tcad)--(metrics);\draw[arr] (metrics)--(score);
\end{tikzpicture}}
\newcommand{\paramtable}{%
\begin{tabularx}{\linewidth}{@{}p{3.6cm}X@{}}\toprule
参数&设定\\\midrule
栅长、栅氧厚度&$L=0.18/0.25/0.35/0.45\um$；主矩阵$5\,\mathrm{nm}$\\
硅体尺寸&宽$L+0.6\um$，深$1.2\um$，有效第三维厚度$1\um$\\
体区、源漏掺杂&P型$10^{17}\,\mathrm{cm}^{-3}$；源漏高斯峰值$10^{20}\,\mathrm{cm}^{-3}$\\
源漏深度参数&$0.12\um$处指定$10^{17}\,\mathrm{cm}^{-3}$，横向因子0.4\\
栅功函数、温度&$4.6\,\mathrm{eV}$；$300\,\mathrm K$\\
电路&存储电容$10\,\mathrm{fF}$；位线与参考电容$100\,\mathrm{fF}$\\
初态&信号SN为$0.1/0.9\V$；中性参考SN为$0.5\V$\\
预充、读出&预充通路$50\,\Omega$，关断$10^{20}\,\Omega$；字线终值$1.8\V$\\\bottomrule
\end{tabularx}}
\newcommand{\timingtable}{%
\begin{tabularx}{\linewidth}{@{}lXlX@{}}\toprule
时间/ns&操作&时间/ns&操作\\\midrule
0.6&释放初态钳位&2.8&提取保持段电荷\\
0.95&检查初态电压&2.8--2.9&位线预充\\
1.0&粒子入射&3.0&断开预充\\
&&3.05--3.10&字线上升，4.0读出\\\bottomrule
\end{tabularx}}
\newcommand{\metricdefs}{%
\begin{align}
\Delta Q&=C_s\,[V_{{\rm SN},e}(2.8\ns)-V_{{\rm SN},c}(2.8\ns)],\\
D_e&=V_{{\rm BL},e}(4\ns)-V_{{\rm BL},n}(4\ns),\qquad M_e=(2b-1)D_e,\\
\Delta M&=M_c-M_e.
\end{align}}
\newcommand{\scoretable}{\begin{tabular}{lllr rrr}
\toprule
集合 & 位 & 目标 & 基准RMSE & 岭回归MAE & RMSE & 最大\\
\midrule
留出 & 0 & $\Delta Q$/fC & 1.9560 & 0.0566 & 0.0883 & 0.1816\\
留出 & 0 & $\Delta M$/mV & 18.4480 & 0.5270 & 0.7984 & 1.6456\\
留出 & 1 & $\Delta Q$/fC & 2.1549 & 0.0618 & 0.0946 & 0.1989\\
留出 & 1 & $\Delta M$/mV & 20.0666 & 0.5498 & 0.8380 & 1.7654\\
外推 & 0 & $\Delta Q$/fC & 2.0254 & 0.0528 & 0.0680 & 0.1364\\
外推 & 0 & $\Delta M$/mV & 18.8677 & 0.4050 & 0.5194 & 1.0753\\
外推 & 1 & $\Delta Q$/fC & 2.1983 & 0.0517 & 0.0672 & 0.1442\\
外推 & 1 & $\Delta M$/mV & 20.2199 & 0.3877 & 0.5049 & 1.1051\\
\bottomrule\end{tabular}
}
